\chapter{Lumbar spine}

\section*{Definition and Overview}
The lumbar spine is the lower region of the vertebral column, typically consisting of five large vertebrae designated L1 to L5 (and occasionally a transitional L6). Situated between the rigid thoracic spine above and the fused sacrum below, the lumbar spine is anatomically designed to bear the majority of the body's axial weight, protect the lower spinal cord and cauda equina, and permit a wide range of movements, including flexion, extension, lateral bending, and rotation. Because of its weight-bearing load and extensive mobility, the lumbar spine is highly susceptible to degenerative changes, mechanical strains, and traumatic injuries, making it the most common source of low back pain.

\section*{Epidemiology}
Low back pain (LBP) is one of the most common medical complaints globally and is the leading cause of years lived with disability worldwide. In Australia, it is estimated that approximately 16\% of the population experiences low back pain at any given time, and up to 80\% of individuals will suffer from it at some point during their lives. The prevalence of lumbar spine disorders peaks between the ages of 40 and 65 years. Degenerative lumbar conditions, such as spinal stenosis and facet joint arthropathy, show a strong positive correlation with advancing age, representing a substantial socioeconomic and healthcare burden.

\section*{Aetiology and Pathophysiology}
Pathology of the lumbar spine arises from structural degradation, mechanical overload, or inflammatory processes affecting the vertebrae, intervertebral discs, joints, ligaments, or muscles. Causes and pathological mechanisms include:

\begin{itemize}
    \item \textbf{Degenerative disc disease:} Age-related dehydration of the nucleus pulposus, which reduces its shock-absorbing capacity, leads to disc height loss and puts mechanical stress on the surrounding annulus fibrosus.
    \item \textbf{Facet joint arthropathy:} Osteoarthritic degeneration of the synovial zygapophyseal (facet) joints, characterised by cartilage erosion, joint space narrowing, and hypertrophic osteophyte formation.
    \item \textbf{Herniated intervertebral disc:} Protrusion or extrusion of the nucleus pulposus through a tear in the annulus fibrosus, causing mechanical compression and chemical irritation of adjacent spinal nerve roots.
    \item \textbf{Lumbar spinal stenosis:} Narrowing of the central spinal canal, lateral recesses, or neural foramina, compressing the cauda equina or spinal nerves; this is typically driven by a combination of facet hypertrophy, osteophytes, and ligamentum flavum thickening.
    \item \textbf{Spondylolisthesis:} The forward displacement of one vertebra relative to the one below it, which can be degenerative or secondary to a defect in the pars interarticularis (spondylolysis).
    \item \textbf{Musculoskeletal strain:} Micro-tears in the paraspinal muscles (such as the erector spinae) or sprains of the lumbar ligaments caused by acute overloading, repetitive twisting, or poor ergonomics.
\end{itemize}

\section*{Clinical Features}
The presentation of lumbar spine disorders depends on the structures involved and the presence of nerve root irritation. Clinical features include:

\begin{itemize}
    \item \textbf{Localized low back pain:} A dull, aching pain, stiffness, and restricted range of motion in the lumbar region, often aggravated by prolonged sitting, standing, or bending.
    \item \textbf{Radicular pain (sciatica):} Sharp, shooting, or electric shock-like pain radiating from the lower back through the buttock and down the posterior or lateral aspect of the leg, extending below the knee in a specific dermatomal distribution (most commonly L4, L5, or S1).
    \item \textbf{Paresthesia and numbness:} Tingling or loss of sensation in the lower limb matching the compressed nerve pathway.
    \item \textbf{Neurogenic claudication:} Bilateral or unilateral leg cramping, pain, or heaviness that is triggered by walking or standing and relieved by sitting or bending forward (which increases the diameter of the spinal canal).
    \item \textbf{Motor deficits and reflex changes:} Weakness in specific muscle groups, such as the tibialis anterior (causing foot drop in L5 pathology) or gastrocnemius (affecting heel-raise in S1 pathology), along with diminished deep tendon reflexes (patellar or Achilles).
\end{itemize}

\section*{Differential Diagnosis}
When evaluating lumbar spine pain, it is vital to distinguish primary mechanical or degenerative spine conditions from inflammatory, vascular, and visceral pathologies:

\begin{itemize}
    \item \textbf{Mechanical back pain vs.\ inflammatory spondyloarthropathy:} Mechanical pain worsens with activity and improves with rest, whereas inflammatory conditions (e.g.\ ankylosing spondylitis) feature morning stiffness lasting over 30 minutes, pain that improves with exercise, and onset in early adulthood.
    \item \textbf{Lumbar radiculopathy vs.\ peripheral vascular disease:} Sciatic nerve pain vs.\ vascular claudication, where cramping leg pain is strictly exertional, relieved quickly by standing still, and associated with cold skin and weak peripheral pulses.
    \item \textbf{Spinal pain vs.\ referred visceral pain:} Pain triggered by movement or spinal palpation vs.\ pain referred from internal organs, such as a leaking abdominal aortic aneurysm, nephrolithiasis (kidney stones), pyelonephritis, or pancreatitis.
    \item \textbf{Degenerative stenosis vs.\ spinal cord neoplasm:} Slow, chronic claudication vs.\ persistent, progressive back pain that is worse when lying flat, disturbs sleep, and is associated with systemic symptoms (e.g.\ unexplained weight loss).
\end{itemize}

\section*{When to Seek Medical Care}
Most episodes of low back pain are benign and self-limiting. However, certain symptoms indicate severe neurological compromise or a life-threatening emergency.

\textbf{Call 000 (or local emergency services) for:}
\begin{itemize}
    \item Sudden loss of bowel or bladder control (urinary retention or fecal incontinence) associated with numbness in the groin, perineum, or buttocks (saddle anaesthesia), indicating cauda equina syndrome.
    \item Sudden, rapidly progressive weakness or paralysis in one or both legs.
    \item Severe, acute back pain accompanied by a tearing or ripping sensation in the abdomen or chest, suggesting a ruptured abdominal aortic aneurysm or aortic dissection.
    \item Severe back pain that occurs immediately after a high-impact fall, motor vehicle accident, or major physical trauma.
    \item Back pain accompanied by high fever, chills, and night sweats, indicating a potential spinal infection (e.g.\ discitis or epidural abscess).
\end{itemize}

\section*{Diagnosis and Investigations}
A thorough history and clinical examination are the cornerstones of diagnosis. Imaging is indicated for red flag symptoms or pain persisting beyond 4 to 6 weeks:

\begin{itemize}
    \item \textbf{Neurological examination:} Assessing gait, spinal range of movement, muscle strength, dermatomal sensation, deep tendon reflexes, and performing nerve stretch tests (such as the Straight Leg Raise or Femoral Stretch tests).
    \item \textbf{Magnetic resonance imaging (MRI):} The gold standard for evaluating soft tissues; it provides high-resolution images of intervertebral discs, nerve roots, ligamentum flavum, and the spinal cord.
    \item \textbf{Plain X-rays:} Used to assess bony alignment, detect spondylolisthesis, evaluate degenerative disc narrowing, and rule out fractures.
    \item \textbf{Computed tomography (CT) scan:} Indicated to evaluate complex bony anatomy, fractures, or in patients who have contraindications to MRI (such as non-compatible pacemakers).
    \item \textbf{Laboratory tests:} Inflammatory markers (ESR, CRP) and blood cultures if spinal infection is suspected, or HLA-B27 typing if an inflammatory spondyloarthropathy is suspected.
\end{itemize}

\section*{Management}
Management follows a stepped care model, starting with conservative therapy and progressing to invasive options if clinical recovery is not achieved:

\begin{itemize}
    \item \textbf{First-line pharmacological therapy:} Short-term use of paracetamol, non-steroidal anti-inflammatory drugs (NSAIDs) to reduce inflammation, and muscle relaxants for acute muscle spasms.
    \item \textbf{Active physical therapy:} Undergoing a structured exercise programme focusing on core stabilisation, lumbar range of motion, hamstring stretching, and low-impact aerobic conditioning (e.g.\ walking or swimming).
    \item \textbf{Interventional pain procedures:} Performing epidural corticosteroid injections for radiculopathy, facet joint injections, or radiofrequency ablation of medial branches for facet arthropathy.
    \item \textbf{Decompressive surgery:} Indicated for progressive neurological deficits or severe pain refractory to conservative therapy; procedures include microdiscectomy, laminectomy, or spinal fusion.
    \item \textbf{Psychological and behavioural support:} Utilizing cognitive behavioural therapy (CBT) and multidisciplinary pain management programmes to improve functional outcomes in chronic pain.
\end{itemize}

\section*{Complications}
Neglecting or delayed treatment of specific lumbar spine pathologies can lead to permanent disability:

\begin{itemize}
    \item \textbf{Chronic pain syndrome and disability:} Chronic, intractable pain leading to loss of mobility, reduced work capacity, depression, and anxiety.
    \item \textbf{Permanent neurological deficit:} Irreversible nerve damage presenting as permanent numbness, muscle wasting, or foot drop.
    \item \textbf{Cauda equina damage:} Chronic bladder or bowel incontinence and sexual dysfunction if emergency decompression is delayed.
    \item \textbf{Postoperative complications:} Surgical risks including dural tear, surgical site infection, failed back surgery syndrome, or adjacent segment disease.
\end{itemize}

\section*{Prognosis and Outcomes}
The prognosis for acute low back pain is excellent, with approximately 90\% of patients experiencing recovery within 6 weeks, regardless of the treatment modality. However, recurrence is common, with up to 50\% of patients experiencing another episode within a year. Approximately 5--10\% of acute cases transition into chronic low back pain, which is more difficult to treat. For patients with progressive radiculopathy or stenosis undergoing appropriate decompression surgery, the outcomes for leg symptoms are highly favourable.

\section*{Prevention and Self-Care}
Engaging in self-care and ergonomic modifications is essential to prevent the onset or recurrence of lumbar spine disorders:

\begin{itemize}
    \item \textbf{Regular core exercises:} Actively engaging in Pilates, yoga, or targeted core-strengthening exercises to provide intrinsic muscular support to the lumbar spine.
    \item \textbf{Ergonomic lift training:} Always bending at the knees rather than the waist, keeping the load close to the body, and avoiding twisting the torso while carrying heavy objects.
    \item \textbf{Maintain a healthy weight:} Minimising excess body weight, particularly abdominal adiposity, to reduce mechanical shear forces on the lumbar vertebrae.
    \item \textbf{Workstation ergonomics:} Using chairs with lumbar support, adjusting desk heights, and taking short movement breaks every 30--45 minutes.
    \item \textbf{Smoking cessation:} Quitting smoking is vital, as nicotine impairs microvascular circulation, accelerating degenerative disc disease and delaying post-surgical healing.
\end{itemize}

\section*{Sources and Further Reading}
The following organisations provide reputable details and guidelines on lumbar spine care:

\begin{itemize}
    \item Healthdirect Australia. \textit{Patient resource on lower back pain causes, symptoms, and self-care tips.} https://www.healthdirect.gov.au/
    \item Mayo Clinic. \textit{Clinical overview of lumbar disc disease, spinal stenosis, and treatment options.} https://www.mayoclinic.org/
    \item NHS. \textit{Detailed guide to back pain, explaining exercises, posture tips, and when to see a GP.} https://www.nhs.uk/
    \item Arthritis Australia. \textit{Consumer information sheet regarding back pain and degenerative joint diseases.} https://arthritisaustralia.com.au/
    \item North American Spine Society. \textit{Evidence-based guidelines and clinical resources for lumbar spine conditions.} https://www.spine.org/
\end{itemize}
